\chapter{Designens eigen logikk}

\begin{motto}
Design er ikkje fråvær av resonnement.\\ Det er ei slutningsform vitskapen ikkje har bruk for.
\end{motto}

Til no har forsvaret vist at design har eit eige studieobjekt og ei eiga kunnskapsform, og at denne kunnskapen er ekte og overførbar. Men ein hardnakka kritikar kan gje seg på alt dette og likevel halde fast ved det han meiner er kjernen: at design manglar ein eigen \emph{logikk}. Vitskapen har deduksjon og induksjon, slutningsformer som lèt seg skrive opp, lære og prøve. Design, vil kritikaren seie, har berre intuisjon der logikken skulle ha vore — eit sprang frå problem til løysing som ingen kan gjere greie for. Dette kapitlet bestrider det. Påstanden er at designresonnementet har ei presis, formell form, at denne forma er studert og lærbar, og at ho er noko vitskapen verken har eller treng. Designtenkinga er ikkje eit hól der ein logikk skulle vore. Ho er ein teori om ein eigen slutningsform.

\section{Abduksjon: den tredje slutningsforma}

Det skarpaste arbeidet her er gjort av Kees Dorst, som tek frasen «designerly thinking» og gjer henne om frå slagord til ein presis logisk struktur \autocite{dorst2011core}. Utgangspunktet hans er det klassiske skiljet mellom slutningsformer. I deduksjon kjenner vi eit prinsipp og eit tilfelle, og sluttar oss til resultatet: gjeve lova og årsaka, finn verknaden. I induksjon kjenner vi tilfelle og resultat, og sluttar oss til prinsippet: gjeve mange årsaker og verknader, finn lova. Begge desse er analytiske; dei rører seg innanfor det som alt finst. Men det finst ein tredje form, som Charles Sanders Peirce kalla \emph{abduksjon}: vi kjenner eit resultat vi vil oppnå og sluttar oss bakover til kva som kunne frambringe det. Det er forklaringa sin logikk, og — det Dorst viser — det er skapinga sin logikk.

Dorst skil to slag abduksjon, og det er her argumentet blir kraftfullt. I den fyrste, vanlege forma kjenner vi både det vi vil oppnå (verdien) og prinsippet som styrer feltet (det «korleis» som gjeld), og vi sluttar oss til kva ting eller grep som vil verke. Dette er ingeniørens og problemløysarens abduksjon. Men i den andre forma — den Dorst kallar kjernen i designtenkinga — kjenner vi berre kva vi vil oppnå. Både \emph{kva} vi skal lage og \emph{etter kva prinsipp} det skal verke, er ukjent. Designaren må finne to ting samstundes: ein ting å lage og ein måte å tenkje om situasjonen på som gjer at den tingen gjev meining. Dette doble sprang — å setje opp eit nytt prinsipp og ein ny ting i lag, slik at dei stør kvarandre — er produktiv abduksjon, og det er ein slutningsform som er like reell som dei to andre, berre retta mot å bringe noko nytt til verda i staden for å analysere det som finst.

Poenget mot \textsc{Grunnlaust} er at dette er ein \emph{logikk}, ikkje fråvær av ein. Han har ein struktur som lèt seg skrive opp, lære og kjenne att. At han ikkje er deduktiv, gjer han ikkje irrasjonell; deduksjon er ikkje einaste forma for gyldig slutning, og har aldri vore det. Når kritikaren seier at designaren «berre hoppar» frå problem til løysing, har han sett noko verkeleg — spranget finst — men han har feiltolka det. Spranget er ikkje eit fall ut av logikken. Det er abduksjon, den eine slutningsforma som faktisk kan produsere det nye, og som vitskapen difor brukar berre i si famlande, hypotesedannande fase, men som design brukar som sitt daglege arbeid.

\section{Innramming som overførbar metode}

Om abduksjonen var berre eit personleg sprang, ville det vere lite vunne; ein logikk som ingen kan lære vidare, er knapt ein logikk. Men Dorst sitt andre store bidrag er å vise at det doble sprang har ein systematisk kjerne som lèt seg trene: \emph{frame creation}, innramminga \autocite{dorst2015frame}. Ei ramme er ein måte å sjå ein situasjon på — eit forslag om at «dersom vi ser dette problemet som om det var \emph{slik}, så opnar desse handlingane seg som meiningsfulle». Den dugande designaren løyser ikkje problemet slik det blir presentert. Han spør fyrst om problemet er rett ramma, og prøver ofte ei ny ramme som får heile situasjonen til å sjå annleis ut — og som dermed gjer løysingar synlege som var usynlege under den gamle.

Dorst byggjer dette ut til ein eksplisitt metodikk med trinn ein kan følgje: bryt situasjonen ned til kjernen, søk etter dei djupare verdiane og temaa som står på spel, og konstruer så ei ny ramme som koplar desse saman på ein måte som opnar for handling. Det avgjerande er at dette er \emph{overførbart}. Det kan lærast, øvast og brukast av folk som ikkje fann det opp; Dorst og medarbeidarane hans har dokumentert korleis metoden er teken i bruk på problem langt utanfor det tradisjonelle designfeltet, frå kriminalitetsførebygging til organisasjonsendring. Ein metode som lèt seg løfte ut av sin opphavlege kontekst og brukast av andre på nye problem, er per definisjon ikkje privat intuisjon. Det er kodifisert kunnskap. Når \textsc{Grunnlaust} seier at design manglar ein overførbar metode, er innramminga eit konkret moteksempel: ein namngjeven, dokumentert, lærbar framgangsmåte for det som er aller mest karakteristisk for design, nemleg å finne det rette spørsmålet før ein søkjer svaret.

\section{Lawson: kognisjonen er studert og avmystifisert}

Alt dette kunne framleis avvisast som teori om design, ikkje observasjon av det. Difor er Bryan Lawson sitt arbeid uunnverleg. I \emph{How Designers Think} sette Lawson seg føre å gjere noko empirisk: å studere korleis designarar faktisk tenkjer, mellom anna gjennom kontrollerte forsøk der han samanlikna korleis designarar og naturvitarar gjekk laus på same oppgåva \autocite{lawson1980}. Funna hans er presise og tåler å gjentakast. Naturvitarane nærma seg oppgåva problemfokusert: dei prøvde å forstå strukturen i problemet fyrst, systematisk, før dei søkte ei løysing. Designarane nærma seg henne løysingsfokusert: dei genererte tidleg framlegg til løysingar og brukte desse til å \emph{utforske} kva problemet eigentleg var. Dette er ikkje slapp metode mot stram. Det er to ulike, rasjonelle strategiar tilpassa to ulike slag oppgåver — og for designoppgåver, der problemet ikkje er fullt kjent før ein prøver å løyse det, er den løysingsfokuserte strategien den fornuftige.

Undertittelen på Lawson si bok er programmet hans: designprosessen \emph{avmystifisert}. Det er nett kva han gjer. Han tek det som såg ut som ein lukka, intuitiv black box og viser at det inni er ein struktur — gjenkjennelege mønster, strategiar, og kjenneteikn på dugleik som skil den røynde frå novisen. Dette er det rake motsette av påstanden om at designtenking er uutgrunneleg. Ho \emph{er} utgrunna; det finst eit halvt hundreår med empirisk forsking på designkognisjon, frå Lawson og framover, som har kartlagt korleis designarar rammar problem, genererer løysingar, og brukar skisser og prototypar til å tenkje med. At resultatet er ein annan kognitiv stil enn vitarens, gjer ikkje stilen mindre verkeleg eller mindre lærbar. Det gjer han til ein eigen, dokumentert måte å tenkje på.

\section{Det kontingente som eige kunnskapsfelt}

Bak alt dette ligg eit filosofisk poeng som bind kapitlet saman, og som ironisk nok kjem frå den same Herbert Simon som drøymde om ein designvitskap \autocite{simon1969sciences}. Simon sitt varige bidrag var å skilje ut det kunstige som eit eige felt for kunnskap: tinga som er slik dei er fordi nokon laga dei slik, og som kunne vore annleis. Naturvitskapen handlar om det naudsynte — om korleis verda \emph{må} vere, gjeve lovene. Design handlar om det \emph{kontingente} — om korleis verda kunne bli, gjeve våre val. Og kunnskap om det kontingente har ein annan logikk enn kunnskap om det naudsynte, fordi objektet er eit anna. Der vitskapen søkjer det allmenne og naudsynte, må design meistre det partikulære og moglege.

Dette forklarar kvifor abduksjonen, innramminga og den løysingsfokuserte kognisjonen ikkje er manglar målte mot vitskapen, men nett dei reiskapane det kontingente krev. Ein kan ikkje \emph{deduere} seg fram til kva ein bør lage, for det finst inga lov som seier kva som bør finnast. Ein kan berre slutte abduktivt mot det moglege, ramme situasjonen på ein fruktbar måte, og prøve grepet mot røynda. Simon såg ikkje dette heilt klart sjølv — han ville framleis temje det kontingente med optimering og algoritme — men den seinare tradisjonen, frå Cross til Dorst, drog konsekvensen: design er kunnskapen om det kontingente, og difor har det ein eigen logikk med naudsyn. Eg skal ikkje overdrive. Det finst ikkje éin samla, lukka formalisme alle designarar deler, slik det finst ein predikatlogikk; designresonnementet er ein familie av studerte slutningsformer, ikkje eit einaste kalkyle. Men det er meir enn nok til å felle den eigentlege skuldinga. Påstanden var at design ikkje har nokon logikk. Sanninga er at det har ein eigen, og at vi veit ein god del om korleis han verkar.

\vignett

Om design har ein eigen logikk, melder det neste spørsmålet seg av seg sjølv: kan denne logikken bere ein metode streng nok til å produsere kunnskap som hopar seg opp over tid? Det er det den neste delen tek føre seg.

\vidarelesing{\fullcite{dorst2011core}; \fullcite{dorst2015frame}; \fullcite{lawson1980}; \fullcite{cross2006}; \fullcite{simon1969sciences}.}
